\documentclass[autodetect-engine,dvi=dvipdfmx,ja=standard,
               a4j,11pt]{bxjsarticle}

\RequirePackage{geometry}
\geometry{reset,paperwidth=210truemm,paperheight=297truemm}
\geometry{hmargin=.75truein,top=20truemm,bottom=25truemm,footskip=10truemm,headheight=0mm}
%\geometry{showframe} % 本文の"枠"を確認したければ，コメントアウト
\usepackage{graphicx}
\usepackage{fancyvrb}
\usepackage{spverbatim}
\renewcommand{\theFancyVerbLine}{\texttt{\footnotesize{\arabic{FancyVerbLine}:}}}

\title{画像処理実験 第1回}
\author{学生番号: 09B24490\\
        加藤弘都 (KATO, HIROTO)}
\date{\number\year 年\number\month 月\number\day 日}

%%======== 本文 ====================================================%%
\begin{document}
\maketitle
% 目次つきの表紙ページにする場合はコメントを外す
%{\footnotesize \tableofcontents \newpage}

%--------------------------------------------------------------------%
\section{課題1} \label{sec:abstract}

\begin{Verbatim}[numbers=left, xleftmargin=10mm, numbersep=6pt,
                    fontsize=\small, baselinestretch=0.8]
def main():
  s=[ord('A'), ord('B'), ord('C')]
  print(s)
  print(bytes(s).decode())
\end{Verbatim}

%--------------------------------------------------------------------%
\begin{figure}[h]
 \centering
 \includegraphics[scale=.5]{0.jpg}
 \caption{入力画像}
 \label{fig:input}
\end{figure}

%--------------------------------------------------------------------%
\section{数式: 行列，ベクトル}

\newcommand\bA{\mbox{\boldmath$A$}}
\newcommand\bb{\mbox{\boldmath$b$}}
太字$\bA$を使えるように定義した．
$\bA^{-1}\bb$等，使う文字を newcommand で定義して使う．

% \boldsymbol{A} は使えない


%--------------------------------------------------------------------%
\section{感想}

%--------------------------------------------------------------------%
% 参考文献
%   以下は，書き方の例である．実際に，参考にした書籍等を見て書くこと．
%   本文で引用する際は，\cite{book:algodata}などとすればよい．
\begin{thebibliography}{99}
  \bibitem{book:algodata} 平田富雄，アルゴリズムとデータ構造，森北出版，1990.
  \bibitem{book:label2} 著者名，書名，出版社，発行年.
  \bibitem{www:label3} WWWページタイトル，\spverb!https://example.of.too.long.url.jp/you.must.use.spverb/and.insert. a.space.somewhere.to.avoid.overfull.html!，アクセス日.
\end{thebibliography}

%--------------------------------------------------------------------%
\end{document}
